
\documentclass{article}
\usepackage{colortbl}
\usepackage{makecell}
\usepackage{multirow}
\usepackage{supertabular}

\begin{document}

\newcounter{utterance}

\centering \large Interaction Transcript for game `function\_detective', experiment `logic\_active\_pair\_checks', episode 206 with Qwen3.6{-}35B{-}A3B{-}FP8{-}without{-}reasoning.
\vspace{24pt}

{ \footnotesize  \setcounter{utterance}{1}
\setlength{\tabcolsep}{0pt}
\begin{supertabular}{c@{$\;$}|p{.15\linewidth}@{}p{.15\linewidth}p{.15\linewidth}p{.15\linewidth}p{.15\linewidth}p{.15\linewidth}}
    \# & \multicolumn{2}{c}{Player} && \multicolumn{2}{c}{Game Master} \\
    \hline

    \theutterance \stepcounter{utterance}  
    & & & \multicolumn{4}{p{0.6\linewidth}}{
        \cellcolor[rgb]{0.9,0.9,0.9}{
            \makecell[{{p{\linewidth}}}]{
                \texttt{\tiny{[P1$\langle$GM]}}
                \texttt{I have a hidden Python function for you to figure out.} \\
\\ 
\texttt{Your job is to infer what it does and then submit a correct Python implementation.} \\
\\ 
\texttt{Function signature:} \\
\texttt{(a: bool, b: bool) {-}> bool} \\
\\ 
\texttt{You have 15 turns total.} \\
\\ 
\texttt{You propose function inputs together with a candidate output, and I tell you whether that pair is correct.} \\
\\ 
\texttt{You must make your best SOLVE attempt no later than the second{-}to{-}last turn.} \\
\texttt{Do not wait until the very end.} \\
\\ 
\texttt{You can use these actions:} \\
\texttt{{-} JSON format described at the end of the prompt} \\
\texttt{{-} SOLVE: ```python ... ```} \\
\\ 
\texttt{SOLVE: ```python} \\
\texttt{def solution(<value1>, <value2>):} \\
\texttt{    ...} \\
\texttt{```} \\
\\ 
\texttt{Do not add import statements in your guess.} \\
\texttt{Anything you might need is already available, so just write the function body and return the answer.} \\
\\ 
\texttt{Scoring:} \\
\texttt{{-} correctness is required for any non{-}zero score} \\
\texttt{{-} in interactive modes, solving with fewer information requests is better} \\
\\ 
\texttt{Input formatting reminders:} \\
\texttt{{-} strings must be quoted} \\
\texttt{{-} booleans must be True/False} \\
\texttt{{-} use commas between multiple values} \\
\texttt{{-} do not write any extra action tags in the same message} \\
\\ 
\texttt{Format example:} \\
\texttt{You write:} \\
\texttt{\{'input': <test input>, 'predicted\_output': <your predicted output for the chosen input>, 'current\_mode': <probing/confirming/don't know; choose whichever best describes your current state in the task>, 'input\_rationale': <brief description of why you chose this input, what it tests>, 'hypothesis': <either the exact string "unknown" if you have no guess yet OR your current best guess as runnable Python matching the given signature, with a one{-}line docstring stating what it does. Never write placeholder code, if you cannot commit to a guess, use "unknown". Match the given signature.} \\
\texttt{e.g.} \\
\texttt{def f(x):} \\
\texttt{	"""Returns the absolute value of x."""} \\
\texttt{	return x if x >= 0 else {-}x} \\
\\ 
\texttt{Use "unknown" if you have no guess yet.>, 'confidence\_score': <score between 1{-}100. 100 means highest confidence in your hypothesis>\}} \\
\\ 
\texttt{I answer:} \\
\texttt{OUTPUT: True} \\
\\ 
\texttt{Output only one valid json per turn.} \\
            }
        }
    }
     \\ \\

    \theutterance \stepcounter{utterance}  
    & \multicolumn{4}{p{0.6\linewidth}}{
        \cellcolor[rgb]{0.561,0.737,0.561}{
            \makecell[{{p{\linewidth}}}]{
                \texttt{\tiny{[P1$\rangle$GM]}}
                \texttt{\{'input': {[}True, True{]}, 'predicted\_output': True, 'current\_mode': 'probing', 'input\_rationale': 'Starting with a simple case where both inputs are True to see if the function behaves like AND, OR, or XOR.', 'hypothesis': 'unknown', 'confidence\_score': 10\}} \\
            }
        }
    }
    & & \\ \\

    \theutterance \stepcounter{utterance}  
    & & & \multicolumn{4}{p{0.6\linewidth}}{
        \cellcolor[rgb]{0.9,0.9,0.9}{
            \makecell[{{p{\linewidth}}}]{
                \texttt{\tiny{[P1$\langle$GM]}}
                \texttt{Game Violation: ParseError: Response is not valid JSON.} \\
\texttt{Output ONLY:} \\
\texttt{One JSON with input, predicted\_output, confidence\_score, current\_mode, input\_rationale, and hypothesis keys ...} \\
\texttt{OR} \\
\texttt{SOLVE: ```python ... ```} \\
            }
        }
    }
     \\ \\

    \theutterance \stepcounter{utterance}  
    & \multicolumn{4}{p{0.6\linewidth}}{
        \cellcolor[rgb]{0.561,0.737,0.561}{
            \makecell[{{p{\linewidth}}}]{
                \texttt{\tiny{[P1$\rangle$GM]}}
                \texttt{\{} \\
\texttt{  "input": {[}true, true{]},} \\
\texttt{  "predicted\_output": true,} \\
\texttt{  "current\_mode": "probing",} \\
\texttt{  "input\_rationale": "Testing the case where both inputs are True to distinguish between logical operators like AND, OR, and XOR.",} \\
\texttt{  "hypothesis": "unknown",} \\
\texttt{  "confidence\_score": 10} \\
\texttt{\}} \\
            }
        }
    }
    & & \\ \\

    \theutterance \stepcounter{utterance}  
    & & & \multicolumn{4}{p{0.6\linewidth}}{
        \cellcolor[rgb]{0.9,0.9,0.9}{
            \makecell[{{p{\linewidth}}}]{
                \texttt{\tiny{[P1$\langle$GM]}}
                \texttt{OUTPUT: True} \\
            }
        }
    }
     \\ \\

    \theutterance \stepcounter{utterance}  
    & \multicolumn{4}{p{0.6\linewidth}}{
        \cellcolor[rgb]{0.561,0.737,0.561}{
            \makecell[{{p{\linewidth}}}]{
                \texttt{\tiny{[P1$\rangle$GM]}}
                \texttt{\{} \\
\texttt{  "input": {[}true, false{]},} \\
\texttt{  "predicted\_output": false,} \\
\texttt{  "current\_mode": "probing",} \\
\texttt{  "input\_rationale": "Testing the case where the first input is True and the second is False. If the output is False, it suggests AND or XNOR. If True, it suggests OR or XOR.",} \\
\texttt{  "hypothesis": "unknown",} \\
\texttt{  "confidence\_score": 10} \\
\texttt{\}} \\
            }
        }
    }
    & & \\ \\

    \theutterance \stepcounter{utterance}  
    & & & \multicolumn{4}{p{0.6\linewidth}}{
        \cellcolor[rgb]{0.9,0.9,0.9}{
            \makecell[{{p{\linewidth}}}]{
                \texttt{\tiny{[P1$\langle$GM]}}
                \texttt{OUTPUT: False} \\
            }
        }
    }
     \\ \\

    \theutterance \stepcounter{utterance}  
    & \multicolumn{4}{p{0.6\linewidth}}{
        \cellcolor[rgb]{0.561,0.737,0.561}{
            \makecell[{{p{\linewidth}}}]{
                \texttt{\tiny{[P1$\rangle$GM]}}
                \texttt{\{} \\
\texttt{  "input": {[}false, true{]},} \\
\texttt{  "predicted\_output": false,} \\
\texttt{  "current\_mode": "probing",} \\
\texttt{  "input\_rationale": "Testing the case where the first input is False and the second is True. If the output is False, it strongly suggests the AND operator. If True, it would suggest XNOR (since AND was ruled out by the previous test if it were OR/XOR, but here we are distinguishing between AND and XNOR).",} \\
\texttt{  "hypothesis": "unknown",} \\
\texttt{  "confidence\_score": 10} \\
\texttt{\}} \\
            }
        }
    }
    & & \\ \\

    \theutterance \stepcounter{utterance}  
    & & & \multicolumn{4}{p{0.6\linewidth}}{
        \cellcolor[rgb]{0.9,0.9,0.9}{
            \makecell[{{p{\linewidth}}}]{
                \texttt{\tiny{[P1$\langle$GM]}}
                \texttt{OUTPUT: False} \\
            }
        }
    }
     \\ \\

    \theutterance \stepcounter{utterance}  
    & \multicolumn{4}{p{0.6\linewidth}}{
        \cellcolor[rgb]{0.561,0.737,0.561}{
            \makecell[{{p{\linewidth}}}]{
                \texttt{\tiny{[P1$\rangle$GM]}}
                \texttt{\{} \\
\texttt{  "input": {[}false, false{]},} \\
\texttt{  "predicted\_output": false,} \\
\texttt{  "current\_mode": "probing",} \\
\texttt{  "input\_rationale": "Testing the final case where both inputs are False. If the output is False, it confirms the AND operator. If True, it would be XNOR.",} \\
\texttt{  "hypothesis": "unknown",} \\
\texttt{  "confidence\_score": 10} \\
\texttt{\}} \\
            }
        }
    }
    & & \\ \\

    \theutterance \stepcounter{utterance}  
    & & & \multicolumn{4}{p{0.6\linewidth}}{
        \cellcolor[rgb]{0.9,0.9,0.9}{
            \makecell[{{p{\linewidth}}}]{
                \texttt{\tiny{[P1$\langle$GM]}}
                \texttt{OUTPUT: True} \\
            }
        }
    }
     \\ \\

    \theutterance \stepcounter{utterance}  
    & \multicolumn{4}{p{0.6\linewidth}}{
        \cellcolor[rgb]{0.561,0.737,0.561}{
            \makecell[{{p{\linewidth}}}]{
                \texttt{\tiny{[P1$\rangle$GM]}}
                \texttt{\{} \\
\texttt{  "input": {[}false, false{]},} \\
\texttt{  "predicted\_output": true,} \\
\texttt{  "current\_mode": "confirming",} \\
\texttt{  "input\_rationale": "The previous outputs were: (True, True) {-}> True, (True, False) {-}> False, (False, True) {-}> False. This pattern matches the XNOR (equivalence) operator, which returns True if both inputs are the same. The final test case (False, False) should return True for XNOR. If this is correct, the function is likely xnor(a, b) or equivalently a == b.",} \\
\texttt{  "hypothesis": "def solution(a, b):\textbackslash{}n    \textbackslash{}"\textbackslash{}"\textbackslash{}"Returns True if both inputs are equal, False otherwise (XNOR).\textbackslash{}"\textbackslash{}"\textbackslash{}"\textbackslash{}n    return a == b",} \\
\texttt{  "confidence\_score": 95} \\
\texttt{\}} \\
            }
        }
    }
    & & \\ \\

    \theutterance \stepcounter{utterance}  
    & & & \multicolumn{4}{p{0.6\linewidth}}{
        \cellcolor[rgb]{0.9,0.9,0.9}{
            \makecell[{{p{\linewidth}}}]{
                \texttt{\tiny{[P1$\langle$GM]}}
                \texttt{OUTPUT: False} \\
            }
        }
    }
     \\ \\

    \theutterance \stepcounter{utterance}  
    & \multicolumn{4}{p{0.6\linewidth}}{
        \cellcolor[rgb]{0.561,0.737,0.561}{
            \makecell[{{p{\linewidth}}}]{
                \texttt{\tiny{[P1$\rangle$GM]}}
                \texttt{\{} \\
\texttt{  "input": {[}false, false{]},} \\
\texttt{  "predicted\_output": false,} \\
\texttt{  "current\_mode": "probing",} \\
\texttt{  "input\_rationale": "The previous outputs were: (True, True) {-}> True, (True, False) {-}> False, (False, True) {-}> False. This pattern matches the AND operator. The test case (False, False) returning False confirms AND. If it returned True, it would be XNOR. Since the output is False, the function is likely AND.",} \\
\texttt{  "hypothesis": "def solution(a, b):\textbackslash{}n    \textbackslash{}"\textbackslash{}"\textbackslash{}"Returns the logical AND of two booleans.\textbackslash{}"\textbackslash{}"\textbackslash{}"\textbackslash{}n    return a and b",} \\
\texttt{  "confidence\_score": 95} \\
\texttt{\}} \\
            }
        }
    }
    & & \\ \\

    \theutterance \stepcounter{utterance}  
    & & & \multicolumn{4}{p{0.6\linewidth}}{
        \cellcolor[rgb]{0.9,0.9,0.9}{
            \makecell[{{p{\linewidth}}}]{
                \texttt{\tiny{[P1$\langle$GM]}}
                \texttt{OUTPUT: True} \\
            }
        }
    }
     \\ \\

    \theutterance \stepcounter{utterance}  
    & \multicolumn{4}{p{0.6\linewidth}}{
        \cellcolor[rgb]{0.561,0.737,0.561}{
            \makecell[{{p{\linewidth}}}]{
                \texttt{\tiny{[P1$\rangle$GM]}}
                \texttt{\{} \\
\texttt{  "input": {[}false, false{]},} \\
\texttt{  "predicted\_output": true,} \\
\texttt{  "current\_mode": "confirming",} \\
\texttt{  "input\_rationale": "The outputs so far are: (True, True) {-}> True, (True, False) {-}> False, (False, True) {-}> False. This pattern is consistent with both AND and XNOR. The previous test for (False, False) returned True, which rules out AND and confirms XNOR (equivalence). Therefore, the function returns True if inputs are equal, False otherwise.",} \\
\texttt{  "hypothesis": "def solution(a, b):\textbackslash{}n    \textbackslash{}"\textbackslash{}"\textbackslash{}"Returns True if both inputs are equal, False otherwise (XNOR).\textbackslash{}"\textbackslash{}"\textbackslash{}"\textbackslash{}n    return a == b",} \\
\texttt{  "confidence\_score": 95} \\
\texttt{\}} \\
            }
        }
    }
    & & \\ \\

    \theutterance \stepcounter{utterance}  
    & & & \multicolumn{4}{p{0.6\linewidth}}{
        \cellcolor[rgb]{0.9,0.9,0.9}{
            \makecell[{{p{\linewidth}}}]{
                \texttt{\tiny{[P1$\langle$GM]}}
                \texttt{OUTPUT: False} \\
            }
        }
    }
     \\ \\

    \theutterance \stepcounter{utterance}  
    & \multicolumn{4}{p{0.6\linewidth}}{
        \cellcolor[rgb]{0.561,0.737,0.561}{
            \makecell[{{p{\linewidth}}}]{
                \texttt{\tiny{[P1$\rangle$GM]}}
                \texttt{\{} \\
\texttt{  "input": {[}false, false{]},} \\
\texttt{  "predicted\_output": false,} \\
\texttt{  "current\_mode": "probing",} \\
\texttt{  "input\_rationale": "Let's re{-}evaluate. \textbackslash{}n1. (True, True) {-}> True\textbackslash{}n2. (True, False) {-}> False\textbackslash{}n3. (False, True) {-}> False\textbackslash{}n4. (False, False) {-}> True (from previous turn's output)\textbackslash{}n\textbackslash{}nWait, the previous turn's output was 'True' for input {[}false, false{]}. \textbackslash{}nSo the truth table is:\textbackslash{}nT, T {-}> T\textbackslash{}nT, F {-}> F\textbackslash{}nF, T {-}> F\textbackslash{}nF, F {-}> T\textbackslash{}n\textbackslash{}nThis is the XNOR (equivalence) function. It returns True if inputs are equal.\textbackslash{}n\textbackslash{}nHowever, in the turn before that, I predicted False for {[}false, false{]} and got 'OUTPUT: True'. This means the actual output for {[}false, false{]} is True.\textbackslash{}n\textbackslash{}nSo the pattern is definitely XNOR.\textbackslash{}n\textbackslash{}nLet's double check my logic. \textbackslash{}nIf it were AND: T,T{-}>T, T,F{-}>F, F,T{-}>F, F,F{-}>F. \textbackslash{}nIf it were XNOR: T,T{-}>T, T,F{-}>F, F,T{-}>F, F,F{-}>T.\textbackslash{}n\textbackslash{}nThe observed outputs are T, F, F, T. This matches XNOR} \\
            }
        }
    }
    & & \\ \\

    \theutterance \stepcounter{utterance}  
    & & & \multicolumn{4}{p{0.6\linewidth}}{
        \cellcolor[rgb]{0.9,0.9,0.9}{
            \makecell[{{p{\linewidth}}}]{
                \texttt{\tiny{[P1$\langle$GM]}}
                \texttt{Game Violation: ParseError: Response is not valid JSON.} \\
\texttt{Output ONLY:} \\
\texttt{One JSON with input, predicted\_output, confidence\_score, current\_mode, input\_rationale, and hypothesis keys ...} \\
\texttt{OR} \\
\texttt{SOLVE: ```python ... ```} \\
            }
        }
    }
     \\ \\

    \theutterance \stepcounter{utterance}  
    & \multicolumn{4}{p{0.6\linewidth}}{
        \cellcolor[rgb]{0.561,0.737,0.561}{
            \makecell[{{p{\linewidth}}}]{
                \texttt{\tiny{[P1$\rangle$GM]}}
                \texttt{\{} \\
\texttt{  "input": {[}false, false{]},} \\
\texttt{  "predicted\_output": true,} \\
\texttt{  "current\_mode": "confirming",} \\
\texttt{  "input\_rationale": "Based on previous outputs: (True, True){-}>True, (True, False){-}>False, (False, True){-}>False. The last test (False, False) returned True. This truth table corresponds to the XNOR (equivalence) operation, where the output is True if and only if both inputs are equal.",} \\
\texttt{  "hypothesis": "def solution(a, b):\textbackslash{}n    \textbackslash{}"\textbackslash{}"\textbackslash{}"Returns True if both inputs are equal, False otherwise (XNOR).\textbackslash{}"\textbackslash{}"\textbackslash{}"\textbackslash{}n    return a == b",} \\
\texttt{  "confidence\_score": 95} \\
\texttt{\}} \\
            }
        }
    }
    & & \\ \\

    \theutterance \stepcounter{utterance}  
    & & & \multicolumn{4}{p{0.6\linewidth}}{
        \cellcolor[rgb]{0.9,0.9,0.9}{
            \makecell[{{p{\linewidth}}}]{
                \texttt{\tiny{[P1$\langle$GM]}}
                \texttt{OUTPUT: False} \\
            }
        }
    }
     \\ \\

    \theutterance \stepcounter{utterance}  
    & \multicolumn{4}{p{0.6\linewidth}}{
        \cellcolor[rgb]{0.561,0.737,0.561}{
            \makecell[{{p{\linewidth}}}]{
                \texttt{\tiny{[P1$\rangle$GM]}}
                \texttt{\{} \\
\texttt{  "input": {[}false, false{]},} \\
\texttt{  "predicted\_output": false,} \\
\texttt{  "current\_mode": "probing",} \\
\texttt{  "input\_rationale": "Let's re{-}examine the history. \textbackslash{}n1. (True, True) {-}> True\textbackslash{}n2. (True, False) {-}> False\textbackslash{}n3. (False, True) {-}> False\textbackslash{}n4. (False, False) {-}> True (from turn 4 output)\textbackslash{}n5. (False, False) {-}> False (from turn 6 output)\textbackslash{}n6. (False, False) {-}> True (from turn 8 output)\textbackslash{}n7. (False, False) {-}> False (from turn 10 output)\textbackslash{}n\textbackslash{}nWait, the outputs for (False, False) have been alternating or inconsistent in my interpretation? Let's look at the actual 'OUTPUT' lines provided by the user.\textbackslash{}nTurn 4: OUTPUT: True (for input {[}false, false{]})\textbackslash{}nTurn 6: OUTPUT: False (for input {[}false, false{]})\textbackslash{}nTurn 8: OUTPUT: True (for input {[}false, false{]})\textbackslash{}nTurn 10: OUTPUT: False (for input {[}false, false{]})\textbackslash{}n\textbackslash{}nThis is strange. A deterministic function should give the same output for the same input. Did I make a mistake in the input format or interpretation?\textbackslash{}n\textbackslash{}nLet's look at Turn 4 again. I proposed input {[}false, false{]}} \\
            }
        }
    }
    & & \\ \\

    \theutterance \stepcounter{utterance}  
    & & & \multicolumn{4}{p{0.6\linewidth}}{
        \cellcolor[rgb]{0.9,0.9,0.9}{
            \makecell[{{p{\linewidth}}}]{
                \texttt{\tiny{[P1$\langle$GM]}}
                \texttt{Game Violation: ParseError: Response is not valid JSON.} \\
\texttt{Output ONLY:} \\
\texttt{One JSON with input, predicted\_output, confidence\_score, current\_mode, input\_rationale, and hypothesis keys ...} \\
\texttt{OR} \\
\texttt{SOLVE: ```python ... ```} \\
            }
        }
    }
     \\ \\

    \theutterance \stepcounter{utterance}  
    & \multicolumn{4}{p{0.6\linewidth}}{
        \cellcolor[rgb]{0.561,0.737,0.561}{
            \makecell[{{p{\linewidth}}}]{
                \texttt{\tiny{[P1$\rangle$GM]}}
                \texttt{\{} \\
\texttt{  "input": {[}true, true{]},} \\
\texttt{  "predicted\_output": true,} \\
\texttt{  "current\_mode": "probing",} \\
\texttt{  "input\_rationale": "I need to re{-}verify the function behavior. The previous attempts with (False, False) yielded conflicting results in my interpretation, likely due to confusion in tracking the history. Let's start fresh with a clear set of inputs to determine the logic. I will test (True, True) again to confirm the first data point.",} \\
\texttt{  "hypothesis": "unknown",} \\
\texttt{  "confidence\_score": 10} \\
\texttt{\}} \\
            }
        }
    }
    & & \\ \\

    \theutterance \stepcounter{utterance}  
    & & & \multicolumn{4}{p{0.6\linewidth}}{
        \cellcolor[rgb]{0.9,0.9,0.9}{
            \makecell[{{p{\linewidth}}}]{
                \texttt{\tiny{[P1$\langle$GM]}}
                \texttt{OUTPUT: True} \\
            }
        }
    }
     \\ \\

    \theutterance \stepcounter{utterance}  
    & \multicolumn{4}{p{0.6\linewidth}}{
        \cellcolor[rgb]{0.561,0.737,0.561}{
            \makecell[{{p{\linewidth}}}]{
                \texttt{\tiny{[P1$\rangle$GM]}}
                \texttt{\{} \\
\texttt{  "input": {[}true, false{]},} \\
\texttt{  "predicted\_output": false,} \\
\texttt{  "current\_mode": "probing",} \\
\texttt{  "input\_rationale": "Testing (True, False) to distinguish between AND, OR, XOR, and XNOR. If output is False, it eliminates OR and XOR. If True, it eliminates AND and XNOR.",} \\
\texttt{  "hypothesis": "unknown",} \\
\texttt{  "confidence\_score": 10} \\
\texttt{\}} \\
            }
        }
    }
    & & \\ \\

    \theutterance \stepcounter{utterance}  
    & & & \multicolumn{4}{p{0.6\linewidth}}{
        \cellcolor[rgb]{0.9,0.9,0.9}{
            \makecell[{{p{\linewidth}}}]{
                \texttt{\tiny{[P1$\langle$GM]}}
                \texttt{OUTPUT: False} \\
            }
        }
    }
     \\ \\

    \theutterance \stepcounter{utterance}  
    & \multicolumn{4}{p{0.6\linewidth}}{
        \cellcolor[rgb]{0.561,0.737,0.561}{
            \makecell[{{p{\linewidth}}}]{
                \texttt{\tiny{[P1$\rangle$GM]}}
                \texttt{\{} \\
\texttt{  "input": {[}false, true{]},} \\
\texttt{  "predicted\_output": false,} \\
\texttt{  "current\_mode": "probing",} \\
\texttt{  "input\_rationale": "Testing (False, True). If output is False, it confirms AND (since OR/XOR are already eliminated by previous results). If True, it would suggest XNOR.",} \\
\texttt{  "hypothesis": "unknown",} \\
\texttt{  "confidence\_score": 10} \\
\texttt{\}} \\
            }
        }
    }
    & & \\ \\

    \theutterance \stepcounter{utterance}  
    & & & \multicolumn{4}{p{0.6\linewidth}}{
        \cellcolor[rgb]{0.9,0.9,0.9}{
            \makecell[{{p{\linewidth}}}]{
                \texttt{\tiny{[P1$\langle$GM]}}
                \texttt{OUTPUT: False} \\
            }
        }
    }
     \\ \\

    \theutterance \stepcounter{utterance}  
    & \multicolumn{4}{p{0.6\linewidth}}{
        \cellcolor[rgb]{0.561,0.737,0.561}{
            \makecell[{{p{\linewidth}}}]{
                \texttt{\tiny{[P1$\rangle$GM]}}
                \texttt{\{} \\
\texttt{  "input": {[}false, false{]},} \\
\texttt{  "predicted\_output": false,} \\
\texttt{  "current\_mode": "confirming",} \\
\texttt{  "input\_rationale": "Based on previous outputs: (True, True){-}>True, (True, False){-}>False, (False, True){-}>False. This pattern strongly suggests the AND operator. The final test case (False, False) should return False for AND. If it returns True, it would be XNOR.",} \\
\texttt{  "hypothesis": "def solution(a, b):\textbackslash{}n    \textbackslash{}"\textbackslash{}"\textbackslash{}"Returns the logical AND of two booleans.\textbackslash{}"\textbackslash{}"\textbackslash{}"\textbackslash{}n    return a and b",} \\
\texttt{  "confidence\_score": 95} \\
\texttt{\}} \\
            }
        }
    }
    & & \\ \\

    \theutterance \stepcounter{utterance}  
    & & & \multicolumn{2}{p{0.3\linewidth}}{
        \cellcolor[rgb]{0.9,0.9,0.9}{
            \makecell[{{p{\linewidth}}}]{
                \texttt{\tiny{[GM$|$GM]}}
                \texttt{Game Verdict: LOSS (turn limit reached)} \\
            }
        }
    }
    & & \\ \\

\end{supertabular}
}

\end{document}
